% Filename: lista1.tex
% 
% This code is part of 'Solutions for MS550, M\'{e}todos de Matem\'{a}tica Aplicada I, and F520, M\'{e}todos Matem\'{a}ticos da F\'{i}sica I'
% 
% Description: This file corresponds to the solutions of homework sheet 1.
% 
% Created: 07.03.12 04:00:00 PM
% Last Change: 07.07.12 05:30:07 PM
% 
% Authors:
% - Raniere Silva (2012): initial version
% 
% Copyright (c) 2012 Raniere Silva <r.gaia.cs@gmail.com>
% 
% This work is licensed under the Creative Commons Attribution-ShareAlike 3.0 Unported License. To view a copy of this license, visit http://creativecommons.org/licenses/by-sa/3.0/ or send a letter to Creative Commons, 444 Castro Street, Suite 900, Mountain View, California, 94041, USA.
%
% This work is distributed in the hope that it will be useful, but WITHOUT ANY WARRANTY; without even the implied warranty of MERCHANTABILITY or FITNESS FOR A PARTICULAR PURPOSE.
%
\documentclass[a4paper,12pt, leqno, answers]{exam}
% Customiza\c{c}\~{a}o da classe exam
\newcommand{\mycheader}{Lista 1 - Sistemas de coordenadas}
\header{MS550, F520}{\mycheader}{\thepage/\numpages}
\headrule
\footer{Dispon\'{i}vel em \\% Filename: repository.tex
% 
% This code is part of 'Solutions for MS550, M\'{e}todos de Matem\'{a}tica Aplicada I, and F520, M\'{e}todos Matem\'{a}ticos da F\'{i}sica I'
% 
% Description: This file keeps the url of the repository.
% 
% Created: 07.03.12 04:00:00 PM
% Last Change: 30.05.12 04:40:25 PM
% 
% Authors:
% - Raniere Silva (2012): initial version
% 
% Copyright (c) 2012 Raniere Silva <r.gaia.cs@gmail.com>
% 
% This work is licensed under the Creative Commons Attribution-ShareAlike 3.0 Unported License. To view a copy of this license, visit http://creativecommons.org/licenses/by-sa/3.0/ or send a letter to Creative Commons, 444 Castro Street, Suite 900, Mountain View, California, 94041, USA.
%
% This work is distributed in the hope that it will be useful, but WITHOUT ANY WARRANTY; without even the implied warranty of MERCHANTABILITY or FITNESS FOR A PARTICULAR PURPOSE.
%
\url{https://github.com/r-gaia-cs/solucoes_listas_metodos}
}{}{Reportar erros para \\\href{mailto:r.gaia.cs@gmail.com}{r.gaia.cs@gmail.com}
}
\footrule 
\pagestyle{headandfoot}
\renewcommand{\solutiontitle}{\noindent\textbf{Solu\c{c}\~{a}o:}\enspace}
\SolutionEmphasis{\slshape}
\unframedsolutions
\pointname{}

% Filename: paper_size.tex
%
% This code is part of 'Solutions for MS550, M\'{e}todos de Matem\'{a}tica Aplicada I, and F520, M\'{e}todos Matem\'{a}ticos da F\'{i}sica I'
% 
% Description: This file corresponds to the paper size output.
% 
% Created: 07.03.12 04:00:00 PM
% Last Change: 30.05.12 04:40:25 PM
% 
% Authors:
% - Raniere Silva (2012): initial version
% 
% Copyright (c) 2012 Raniere Silva <r.gaia.cs@gmail.com>
% 
% This work is licensed under the Creative Commons Attribution-ShareAlike 3.0 Unported License. To view a copy of this license, visit http://creativecommons.org/licenses/by-sa/3.0/ or send a letter to Creative Commons, 444 Castro Street, Suite 900, Mountain View, California, 94041, USA.
%
% This work is distributed in the hope that it will be useful, but WITHOUT ANY WARRANTY; without even the implied warranty of MERCHANTABILITY or FITNESS FOR A PARTICULAR PURPOSE.
%
% Para impress\~{a}o
\usepackage[top=3cm, bottom=3cm, left=2cm, right=2cm]{geometry}

% Para ereaders (Kindle, Nook, Kobo, ...) and tablets (iPad, GalaxyTab, ...)
% \usepackage[papersize={180mm,240mm},margin=2mm]{geometry}
% \sloppy


% Pacotes
\usepackage[utf8]{inputenc}
\usepackage[T1]{fontenc}
\usepackage[brazil]{babel}
\usepackage{amsmath}
\usepackage{amsfonts}
\usepackage{amssymb}
\usepackage{hyperref}
\usepackage{graphicx}

% Customiza\c{c}\~{a}o do pacote amsmath
\allowdisplaybreaks[4]

% Novos ambientes
% \newenvironment{fwsolution}{\begin{EnvFullwidth}\begin{TheSolution}}{\end{TheSolution}\end{EnvFullwidth}}

% Novos comandos
\newcommand{\devp}[2]{\frac{\partial #1}{\partial #2}}
\newcommand{\grad}{\mbox{grad }}
\newcommand{\diver}{\mbox{div }}
\newcommand{\rot}{\mbox{rot }}

\begin{document}
%cover
\thispagestyle{empty}
% Filename: cover.tex
% This code is part of 'Solutions for MS550, M\'{e}todos de Matem\'{a}tica Aplicada I, and F520, M\'{e}todos Matem\'{a}ticos da F\'{i}sica I'
% 
% Description: This file corresponds to the cover.
% 
% Created: 30.05.12 04:40:25 PM
% Last Change: 31.05.12 10:11:55 PM
% 
% Authors:
% - Raniere Silva (2012): initial version
% 
% Copyright (c) 2012 Raniere Silva <r.gaia.cs@gmail.com>
% 
% This work is licensed under the Creative Commons Attribution-ShareAlike 3.0 Unported License. To view a copy of this license, visit http://creativecommons.org/licenses/by-sa/3.0/ or send a letter to Creative Commons, 444 Castro Street, Suite 900, Mountain View, California, 94041, USA.
%
% This work is distributed in the hope that it will be useful, but WITHOUT ANY WARRANTY; without even the implied warranty of MERCHANTABILITY or FITNESS FOR A PARTICULAR PURPOSE.
%
\begin{center}
    \LARGE{Solu\c{c}\~{o}es para MS550, M\'{e}todos de Matem\'{a}tica Aplicada I, e F520, M\'{e}todos Matem\'{a}ticos da F\'{i}sica I}
    
    \Large{\mycheader}
\end{center}
\vspace{.5\textheight}

\begin{tabular}{|p{.9\textwidth}|}
\hline
Este trabalho foi licenciado com a Licen\c{c}a Creative Commons Atribui\c{c}\~{a}o - CompartilhaIgual 3.0 N\~{a}o Adaptada. Para ver uma c\'{o}pia desta licen\c{c}a, visite \url{http://creativecommons.org/licenses/by-sa/3.0/} ou envie um pedido por carta para Creative Commons, 444 Castro Street, Suite 900, Mountain View, California, 94041, USA.
\begin{center}
\includegraphics[scale=1]{cc-by-sa.png}
\end{center}
Este trabalho \'{e} distribuido na esperança que possa ser \'{u}til, mas SEM NENHUMA GARANTIA; sem uma garantia implicita de ADEQUA\c{C}\~{A}O a qualquer MERCADO ou APLICA\c{C}\~{A}O EM PARTICULAR.
\\ \hline
\end{tabular}

\newpage
\setcounter{page}{1}
Algumas express\~{o}es \'{u}teis na resolu\c{c}\~{a}o das quest\~{o}es:
\begin{align}
    & \devp{\vec{r}}{q_i} = \devp{x}{q_i} \vec{e}_x + \devp{y}{q_i} \vec{e}_y + \devp{z}{q_i} \vec{e}_z,
    \label{eq:vetor_tangente} \\
    & h_i = \sqrt{\left(\devp{x}{q_i}\right)^2 + \left(\devp{y}{q_i}\right)^2 + \left(\devp{z}{q_i}\right)^2},
    \label{eq:fator_escala} \\
    & \vec{e}_{q_i} = \frac{1}{h_i} \devp{\vec{r}}{q_i}
    \label{eq:vetor_tang_unit} \\
    &\vec{e}_{q_i} = \sum_j \left( \vec{e}_{q_i} \cdot \vec{e}_{p_j} \right) \vec{e}_{p_j}
    \label{eq:vetor_mud_coor} \\
    & \grad f = \nabla f = \frac{1}{h_1} \devp{f}{q_1} \vec{e}_{q_1} + \frac{1}{h_2} \devp{f}{q_2} \vec{e}_{q_2} + \frac{1}{h_3} \devp{f}{q_3} \vec{e}_{q_3},
    \label{eq:grad} \\
    & \diver \vec{V} = \nabla \cdot \vec{V} = \frac{1}{h_1 h_2 h_3} \left[\devp{h_2 h_3 V_1}{q_1} + \devp{h_3 h_1 V_2}{q_2} + \devp{h_1 h_2 V_3}{q_3}\right],
    \label{eq:div} \\
    & \rot \vec{V} = \nabla \times \vec{V} = \frac{1}{h_1 h_2 h_3} \begin{vmatrix}
        h_1 \vec{e}_{q_1} & h_2 \vec{e}_{q_2} & h_3 \vec{e}_{q_3} \\
        \devp{}{q_1}      & \devp{}{q_2}      & \devp{}{q_3}      \\
        h_1 V_1           & h_2 V_2           & h_3 V_3           
    \end{vmatrix},
    \label{eq:rot} \\
    \begin{split}
        & \nabla^2 f = \frac{1}{h_1 h_2 h_3} \left[\devp{}{q_1}\left(\frac{h_2 h_3}{h_1} \devp{}{q_1}\right)\right. \\ &\qquad\quad\left. + \devp{}{q_2}\left(\frac{h_1 h_3}{h_2} \devp{}{q_2}\right) + \devp{}{q_3}\left(\frac{h_1 h_2}{h_3} \devp{}{q_3}\right)\right] f,
    \end{split}
    \label{eq:laplaciano} \\
    & \begin{cases}
        x = \rho \cos \theta \\
        y = \rho \sin \theta \\
        z = z
    \end{cases}, && \text{coordenadas cil\'{i}ndricas}
    \label{eq:coor_cil} \\
    & \begin{cases}
        x = r \sin \theta \cos \phi \\
        y = r \sin \theta \sin \phi \\
        z = r \cos \theta
    \end{cases}, && \text{coordenadas esf\'{e}ricas}
    \label{eq:coor_esf} \\
    & \nabla \cdot (f \vec{V}) = \vec{V} \cdot \nabla f + f \nabla \cdot \vec{V}, \\
    & \nabla \times (f \vec{V} = f \nabla \times \vec{V} + \nabla f \times \vec{V}, \\
    & \nabla(f g) = f \nabla g + g \nabla f.
\end{align}

\begin{questions}
    \question Seja $\vec{V} = z\vec{i} - 2x\vec{j} + y \vec{k}$. Mostre que as componentes de $\vec{V}$ em coordenadas cil\'{i}ndricas circulares s\~{a}o dadas por
    \begin{align*}
        V_\rho &= z \cos \theta - 2 \rho \cos \theta \sin \theta, \\
        V_\theta &= -z \sin \theta - 2 \rho \cos^2 \theta, \\
        V_z &= \rho \sin \theta
    \end{align*}
    \begin{solution}
        Usando \eqref{eq:coor_cil} escrevemos $\vec{V}$ como
        \[
        \vec{V} = z\vec{i} - 2 \rho \cos \theta \vec{j} + \rho \sin \theta \vec{k}.
        \]
  
        Calculamos os fatores de escala utilizando \eqref{eq:fator_escala}:
        \begin{align*}
            h_p &= \sqrt{\cos^2 \theta + \sin^2 \theta} = 1, \\
            h_\theta &= \sqrt{\rho^2 \sin^2 \theta + \rho^2 \cos^2 \theta} = \rho, \\
            h_z &= \sqrt{1} = 1
        \end{align*}
        e com eles os vetores os vetores tangentes unit\'{a}rios dados por \eqref{eq:vetor_tang_unit}:
        \begin{align*}
            \vec{e}_\rho &= \cos \theta \vec{j} + \sin \theta \vec{k}, \\
            \vec{e}_\theta &= - \sin \theta \vec{j} + \cos \theta \vec{k}, \\
            \vec{e}_z &= \vec{i}.
        \end{align*}
  
        Utilizando \eqref{eq:vetor_mud_coor} escrevemos $\vec{i}$, $\vec{j}$ e $\vec{k}$ em fun\c{c}\~{a}o de $\vec{e}_\rho$, $\vec{e}_\theta$ e $\vec{e}_z$:
        \begin{align*}
            \vec{i} &= \vec{e}_z, \\
            \vec{j} &= \cos \theta \vec{e}_\rho - \sin \theta \vec{e}_\theta, \\
            \vec{k} &= \sin \theta \vec{e}_\rho + \cos \theta \vec{e}_\theta.
        \end{align*}
  
        Reescrevendo $\vec{V}$ em fun\c{c}\~{a}o dos vetores $\vec{e}_\rho$, $\vec{e}_\theta$ e $\vec{e}_z$ obtemos
        \begin{align*}
            \vec{V} &= z \vec{e}_z - 2 \rho \cos \theta \left(-2 \cos \theta \vec{e}_\rho + 2 \sin \theta \vec{e}_\theta\right) + \rho \sin \theta \left(\sin \theta \vec{e}_\rho + \cos \theta \vec{e}_\theta\right) \\
            &= \left(4 \rho \cos^2 \theta + \rho \sin^2 \theta\right) \vec{e}_\rho + \left( -2 \rho \sin \theta \cos \theta + \rho \sin \theta \cos \theta\right) \vec{e}_\theta + z \vec{e}_z \\
            &= \left(3 \rho \cos^2 \theta + \rho\right) \vec{e}_\rho - \rho \sin \theta \cos \theta \vec{e}_\theta + z \vec{e}_z.
        \end{align*}
    \end{solution}
  
    \question Seja o campo vetorial
    \[
    \vec{V} = V_\rho(\rho, \theta) \vec{e}_\rho + V_\theta(\rho, \theta) \vec{e}_\theta.
    \]
    Mostre que $\nabla \times \vec{V} = \mbox{rot } \vec{V}$ possui componente apenas na dire\c{c}\~{a}o $z$.
    \begin{solution}
        Utilizando \eqref{eq:rot} temos
        \begin{align*}
            \nabla \times \vec{V} &= \begin{vmatrix}
                  \vec{e}_\rho & \vec{e}_\theta & \vec{e}_z \\
                  \devp{}{\rho} & \devp{}{\theta} & \devp{}{z} \\
                  V_\rho & V_\theta & 0
            \end{vmatrix} \\
            &= - \devp{V_\theta}{z} \vec{e}_\rho + \devp{V_\rho}{z} \vec{e}_\theta + \left(\devp{V_\theta}{p} - \devp{V_\rho}{\theta}\right) \vec{e}_z \\
            &= \left(\devp{V_\theta}{p} - \devp{V_\rho}{\theta}\right) \vec{e}_z.
        \end{align*}
    \end{solution}
  
    \question Seja o campo vetorial $\vec{V} = \rho \vec{e}_z$. Mostre que
    \begin{align*}
        \nabla \times \vec{V} &= -\vec{e}_\theta, &
        \nabla \times (\vec{V} \times (\nabla \times \vec{V})) &= 0.
    \end{align*}
    \begin{solution}
        Para mostrar que $\nabla \times \vec{V} = - \vec{e}_\theta$ utilizamos \eqref{eq:rot}:
        \begin{align*}
            \nabla \times \vec{V} &= \begin{vmatrix}
                  \vec{e}_\rho & \vec{e}_\theta & \vec{e}_z \\
                  \devp{}{\rho} & \devp{}{\theta} & \devp{}{z} \\
                  0 & 0 & \rho
            \end{vmatrix} \\
            &= \devp{\rho}{\theta} \vec{e}_\rho - \devp{\rho}{\rho} \vec{e}_\theta \\
            &= - \vec{e}_\theta.
        \end{align*}
  
        E para mostrar $\nabla \times \left(\vec{V} \times \left(\nabla \times \vec{V}\right)\right) = 0$ come\c{c}amos utilizando o resultado anterior, i.e., $\nabla \times \vec{V} = - \vec{e}_\rho$:
        \begin{align*}
            \vec{V} \times \left(\nabla \times \vec{V}\right) &= \vec{V} \times \left(- \vec{e}_\theta\right) \\
            &= \rho \vec{e}_z \times \left(- \vec{e}_\theta\right) \\
            &= 0 \\
            \nabla \times \left(\vec{V} \times \left(\nabla \times \vec{V}\right)\right) &= \nabla \times 0 \\
            &= 0.
        \end{align*}
    \end{solution}
  
    \question Calcule $\devp{\vec{e}_{q_i}}{q_j}$ com $i,j = 1, 2, 3$ quando $q_i$ s\~{a}o as coordenadas cil\'{i}ndricas e mostre que as únicas derivadas n\~{a}o nulas s\~{a}o
    \begin{align*}
        \devp{\vec{e}_\rho}{\theta} &= \vec{e}_\theta, \\
        \devp{\vec{e}_\theta}{\theta} &= - \vec{e}\rho.
    \end{align*}
    \begin{solution}
        Utilizando \eqref{eq:coor_cil} e \eqref{eq:fator_escala} temos
        \begin{align*}
            h_\rho &= \sqrt{\cos^2 \theta + \sin^2 \theta} = 1, \\
            h_\theta &= \sqrt{\rho^2 \sin^2 \theta + \rho \cos^2 \theta} = \rho, \\
            h_z &= \sqrt{1} = 1.
        \end{align*}
  
        Calculando os vetores tangentes unit\'{a}rios temos
        \begin{align*}
            \vec{e}_\rho &= \cos \theta \vec{i} + \sin \theta \vec{j}, \\
            \vec{e}_\theta &= \frac{1}{\rho} \left(-\rho \sin \theta \vec{i} + \rho \cos \theta \vec{j}\right) = -\sin \theta \vec{i} + \cos \theta \vec{j}, \\
            \vec{e}_z &= \vec{k}.
        \end{align*}
  
        Ent\~{a}o
        \begin{align*}
            \devp{\vec{e}_\rho}{\rho} &= 0, & 
            \devp{\vec{e}_\rho}{\theta} &= -\sin \theta \vec{i} + \cos \theta \vec{j} = \vec{e}_\theta, &
            \devp{\vec{e}_\rho}{z} &= 0, \\
            \devp{\vec{e}_\theta}{\rho} &= 0, &
            \devp{\vec{e}_\theta}{\theta} &= - \cos \theta \vec{i} - \sin \theta \vec{j} = - \vec{e}_\rho, &
            \devp{\vec{e}_\theta}{z} &= 0, \\
            \devp{\vec{e}_z}{\rho} &= 0, &
            \devp{\vec{e}_z}{\theta} &= 0, &
            \devp{\vec{e}_z}{z} &= 0.
        \end{align*}
    \end{solution}
  
    \question Sejam as coordenadas esferoidais achatadas $(\xi, \eta, \phi)$ dadas por
    \begin{align*}
        x &= a \cosh \xi \cos \eta \cos \phi, \\
        y &= a \cosh \xi \cos \eta \sin \phi, \\
        z &= a \sinh \eta \sin \eta,
    \end{align*}
    onde $\xi \geq 0$, $-\pi/2 \leq \eta \leq \pi/2$, $0 \leq \phi \leq 2\pi$.
    
    Mostre que os fatores de escala s\~{a}o dados por
    \begin{align*}
        h_\xi = h\eta &= a \sqrt{\sinh^2 \xi + \sin^2 \eta}, \\
        h_\phi &= a \cosh \xi \cos \eta.
    \end{align*}
    \begin{solution}
        Utilizando \eqref{eq:fator_escala} vamos calcular os fatores de escala:
        \begin{align*}
            h_\xi &= \sqrt{\left(a \sinh \xi \cos \eta \cos \phi\right)^2 + \left(a \sinh \xi \cos \eta \sin \phi\right)^2 + \left(a \cosh \xi \sin \eta\right)^2} \\
            &= a \sqrt{\sinh^2 \xi \cos^2 \eta \left(\cos^2 \phi + \sin^2 \phi \right) + \cosh^2 \xi \sin^2 \eta} \\
            &= a \sqrt{\sinh^2 \xi \left(\cos^2 \eta + \sin^2 \eta\right) + \sin^2 \eta} \\
            &= a \sqrt{\sinh^2 \xi + \sin^2 \eta}, \\
            h_\eta &= \sqrt{\left(a \cosh \xi \sin \eta \cos \phi\right)^2 + \left(a \cosh \xi \sin \eta \sin \phi\right)^2 + \left(a \sinh \xi \cos \eta\right)^2} \\
            &= a \sqrt{\cosh^2 \xi \sin^2 \eta \left(\cos^2 \phi + \sin^2 \phi\right) + \sinh^2 \xi \cos^2 \eta} \\
            &= a \sqrt{\left(1 + \sinh^2 \xi\right) \sin^2 \eta + \sinh^2 \cos^2 \eta} \\
            &= a \sqrt{\sin^2 \eta + \sinh^2 \xi}, \\
            h_\phi &= \sqrt{\left(a \cosh \xi \cos \eta \sin \phi\right)^2 + \left(a \cosh \xi \cos \eta \cos \phi\right)^2 + 0} \\
            &= a \sqrt{\cosh^2 \xi \cos^2 \eta \left(\sin^2 \phi + \cos^2 \phi\right)} \\
            &= a \cosh \xi \cos \eta.
        \end{align*}
    \end{solution}
  
    \question Sejam $(u, v, z)$ as coordenadas cil\'{i}ndricas parab\'{o}licas, definidas como
    \[
    x = \frac{1}{2} (u^2 - v^2), \ y = uv, \ z = z,
    \]
    com $-\infty < u < +\infty$, $v \geq 0$, $-\infty < z < +\infty$, e $\vec{r}$ o vetor posi\c{c}\~{a}o, $\vec{r} = x \vec{i} + y \vec{j} + z \vec{k}$. Mostre que em coordenadas cil\'{i}ndricas parab\'{o}licas temos
    \[
    \vec{r} = \frac{1}{2} \sqrt{u^2 + v^2} (u \vec{e}_u + v \vec{e}_v) + z \vec{e}_z.
    \]
    Usando essas coordenadas, mostre que $\mbox{div } \vec{r} = 3$.
    \begin{solution}
        Sendo
        \[
        \vec{r} = \frac{1}{2} \left(u^2 - v^2\right) \vec{i} + u v \vec{j} + z \vec{k}
        \]
        come\c{c}amos utilizando \eqref{eq:fator_escala} para calcular os fatores de escala:
        \begin{align*}
            h_u &= \sqrt{u^2 + v^2}, \\
            h_v &= \sqrt{u^2 + v^2}, \\
            h_z &= 1.
        \end{align*}
        Em seguida utilizamos \eqref{eq:vetor_tangente} para calcular os vetores tangentes:
        \begin{align*}
            \vec{e}_u &= \frac{1}{\sqrt{u^2 + v^2}} \left(u \vec{i} + v \vec{j}\right), \\
            \vec{e}_v &= \frac{1}{\sqrt{u^2 + v^2}} \left(-v \vec{i} + u \vec{j}\right), \\
            \vec{e}_z &= \vec{k}.
        \end{align*}
  
        Utilizando \eqref{eq:vetor_mud_coor}, escrevemos $\vec{i}$, $\vec{j}$, $\vec{k}$ em fun\c{c}\~{a}o de $\vec{u}$, $\vec{v}$ e $\vec{z}$:
        \begin{align*}
            \vec{i} &= \frac{1}{\sqrt{u^2 + v^2}} \left(u \vec{e}_u - v \vec{e}_v\right), \\
            \vec{j} &= \frac{1}{\sqrt{u^2 + v^2}} \left(v \vec{e}_u + u \vec{e}_v\right), \\
            \vec{k} &= \vec{e}_z.
        \end{align*}
  
        Substituindo $\vec{i}$, $\vec{j}$ e $\vec{k}$ em $\vec{r}$ temos
        \begin{align*}
            \vec{r} &= \frac{u^2 - v^2}{2 \sqrt{u^2 + v^2}} \left( u \vec{e}_u - v \vec{e}_v \right) + \frac{u v}{\sqrt{u^2 + v^2}} \left( v \vec{e}_u + u \vec{e}_v \right) + z \vec{e}_z \\
            &= \frac{1}{2 \sqrt{u^2 + v^2}} \left( u^3 \vec{e}_u - u^2 v \vec{e}_v - u v^2 \vec{e}_u + v^3 \vec{e}_v + 2 u v^2 \vec{e}_u + 2 u^2 v \vec{e}_v \right) + z \vec{e}_z \\
            &= \frac{1}{2 \sqrt{u^2 + v^2}} \left( u^3 \vec{e}_u + u^2 v \vec{e}_v + u v^2 \vec{e}_u + v^3 \vec{e}_v \right) + z \vec{e}_z \\
            &= \frac{1}{2 \sqrt{u^2 + v^2}} \left( u^2 \left( u \vec{e}_u + v \vec{e}_v \right) + v^2 \left( u \vec{e}_u + v \vec{e}_v \right) \right) + z \vec{e}_z \\
            &= \frac{u^2 + v^2}{2 \sqrt{u^2 + v^2}} \left( u \vec{e}_u + v \vec{e}_v \right) + z \vec{e}_z \\
            &= \frac{\sqrt{u^2 + v^2}}{2} \left( u \vec{e}_u + v \vec{e}_v \right) + z \vec{e}_z
        \end{align*}
        
        Por \'{u}ltimo, utilizando \eqref{eq:div} temos
        \begin{align*}
            \diver \vec{r} &= \frac{1}{h_u h_v h_z} \left[\devp{h_v h_z V_u}{u} + \devp{h_z h_u V_v}{v} + \devp{h_u h_v V_z}{q_z}\right] \\
            &= \frac{1}{u^2 + v^2} \left( \devp{\left( u^2 + v^2 \right) u \ 2}{u} + \devp{\left( u^2 + v^2 \right) v \ 2}{v} + \devp{\left( u^2 + v^2 \right) z}{z} \right) \\
            &= \frac{1}{u^2 + v^2} \left( \frac{2 u^2 + u^2 + v^2}{2} + \frac{2 v^2 + u^2 + v^2}{2} + \left( u^2 + v^2 \right) \right) \\
            &= \frac{1}{u^2 + v^2} \left( \frac{6 \left( u^2 + v^2 \right)}{2} \right) \\
            &= 3
        \end{align*}
    \end{solution}
  
    \question Sejam $(u, v, z)$ as coordenadas cil\'{i}ndricas parab\'{o}licas e o campo vetorial
    \[
    \vec{V} = \frac{1}{4} \sqrt{u^2 + v^2} (-v \vec{e}_u + u \vec{e}_v).
    \]
    Mostre que $\rot \vec{V} = \vec{e}_z$.
    \begin{solution}
        Para utilizar \eqref{eq:rot} precisamos do fator de escala e das componentes dos vetores tangentes. Primeira vamos calcular os fatores de escala utilizando \eqref{eq:fator_escala}:

        \begin{align*}
              h_u &= \sqrt{u^2 + v^2}, \\
              h_v &= \sqrt{v^2 + u^2}, \\
              h_z &= 1.
        \end{align*}
  
        Para calcular as componentes dos vetores tangentes note que
        \[
        V_i = \vec{V} \vec{e}_i.
        \]
        Logo,
        \begin{align*}
            V_u &= -\frac{v}{4} \sqrt{u^2 + v^2}, \\
            V_v &= \frac{u}{4} \sqrt{u^2 + v^2}, \\
            V_z &= 0.
        \end{align*}
         
        Por fim, calculando $\rot \vec{V}$ temos
        \begin{align*}
            \rot \vec{V} &= \frac{1}{u^2 + v^2} \begin{vmatrix}
                \sqrt{u^2 + v^2} \vec{e}_{u} & \sqrt{u^2 + v^2} \vec{e}_{v} & 1 \vec{e}_{z} \\
                \devp{}{u}      & \devp{}{v}      & \devp{}{z}      \\
                -\frac{v}{4} \left(u^2 + v^2)\right) & \frac{u}{4} \left(u^2 + v^2)\right) & 0           
            \end{vmatrix} \\
            &= \frac{1}{u^2 + v^2} \left(\vec{e}_z \left(\frac{1}{4}\left(3u^2 + v^2\right)\right) - \vec{e}_z \left(-\frac{1}{4}\left(u^2 + 3v^2\right)\right)\right) \\
            &= \frac{1}{u^2 + v^2}\left(\frac{\vec{e}_z}{4}\left(3u^2 + v^2 + u^2 + 3v^2\right)\right) \\
            &= \frac{1}{u^2 + v^2}\left(\frac{\vec{e}_z}{4}\left(4u^2 + 4v^2\right)\right) \\
            &= \vec{e}_z.
        \end{align*}
    \end{solution}
  
    \question Sejam as coordenadas $(u, v, z)$ definidas como
    \[
    x = \frac{1}{2} (u^2 + v^2), \ y = uv, \ z = z.
    \]
    Mostre que esse sistema de coordenadas n\~{a}o \'{e} ortogonal.
    \begin{solution}
        Mostrar que esse sistema de coordendas n\~{a}o \'{e} ortogonal equivale a mostrar que existe pelo menos um par $(i,j)$, $i, j = 1, 2, 3$, tal que $\vec{e}_{q_i} \cdot \vec{e}_{q_j} \neq 0$.
  
        Come\c{c}amos calculando os vetores tangentes
        \begin{align*}
            \devp{\vec{r}}{u} &= u \vec{i} + v \vec{j}, \\
            \devp{\vec{r}}{v} &= v \vec{i} + u \vec{j}, \\
            \devp{\vec{r}}{z} &= \vec{k}.
        \end{align*}
        Agora tentamos descobrir um par $(i,j)$, $i, j = 1, 2, 3$, tal que $\vec{e}_{q_i} \cdot \vec{e}_{q_j} \neq 0$:
        \begin{align*}
            \devp{\vec{r}}{u} \cdot \devp{\vec{r}}{v} &= u v + u v = 2 uv \neq 0.
        \end{align*}
    \end{solution}
  
    \question Calcule $\devp{\vec{e}_{q_i}}{q_j}$ com $i, j = 1, 2, 3$ quando $q_i$ s\~{a}o as coordenadas esf\'{e}ricas e mostre que as únicas derivadas n\~{a}o nulas s\~{a}o
    \begin{align*}
        \devp{\vec{e}_r}{\phi} &= \sin \theta \vec{e}_\phi, &
        \devp{\vec{e}_\phi}{\phi} &= -\cos \theta \vec{e}_\theta - \sin \theta \vec{e}_r, \\
        \devp{\vec{e}_\theta}{\phi} &= \cos \theta \vec{e}_\phi, &
        \devp{\vec{e}_r}{\theta} &= \vec{e}_\theta, &
        \devp{\vec{e}_\theta}{\theta} &= -\vec{e}_r.
    \end{align*}
    \begin{solution}
        Seja $\vec{r} = x \vec{i} + y \vec{j} + z \vec{k}$ e $x$, $y$ e $z$ dados por \eqref{eq:coor_esf}. Primeiramente vamos calcular os fatores de escala dados por \eqref{eq:fator_escala}:
        \begin{align*}
            h_r &= \sqrt{\sin^2 \theta \cos^2 \phi + \sin^2 \theta \sin^2 \phi + \cos^2 \theta} = \sqrt{\sin^2  + \cos^2 \theta} = 1, \\
            h_\theta &= \sqrt{r^2 \cos^2 \theta \cos^2 \phi + r^2 \cos^2 \theta \sin \phi + r^2 \sin^2 \theta} = \sqrt{r} = r, \\
            h_\phi &= \sqrt{r^2 \sin^2 \theta \sin^2 \phi + r^2 \sin^2 \theta \cos^2 \phi} = r \sin \theta.
        \end{align*}
  
        Depois calculamos os vetores tangentes unit\'{a}rios dados por \eqref{eq:vetor_tang_unit}:
        \begin{align*}
            \vec{e}_r &= \sin \theta \cos \phi \vec{i} + \sin \theta \sin \phi \vec{j} + \cos \theta \vec{k}, \\
            \vec{e}_\theta &= \cos \theta \cos \phi \vec{i} + \cos \theta \sin \phi \vec{j} - \sin \theta \vec{k}, \\
            \vec{e}_\phi &= - \sin \phi \vec{i} + \cos \phi \vec{j}.
        \end{align*}
  
        Agora podemos calcular as derivadas dos vetores tangentes unit\'{a}rios:
        \begin{align*}
            \devp{\vec{e}_r}{\theta} &= \cos \theta \cos \phi \vec{i} + \cos \theta \sin \phi \vec{j} - \sin \theta \vec{k} = \vec{e}_\theta, &
            \devp{\vec{e}_r}{r} &= 0, \\
            \devp{\vec{e}_r}{\phi} &= -\sin \theta \sin \phi \vec{i} + \sin \theta \cos \phi \vec{j} = \sin \theta \vec{e}_\phi, &
            \devp{\vec{e}_\theta}{r} &= 0 \\
            \devp{\vec{e}_\theta}{\theta} &= - \sin \theta \cos \phi \vec{i} - \sin \theta \sin \phi \vec{j} - \cos \theta \vec{k} = -\vec{e}_r, &
            \devp{\vec{e}_\phi}{r} &= 0, \\
            \devp{\vec{e}_\theta}{\phi} &= - \cos \theta \sin \phi \vec{i} + \cos \theta \cos \phi \vec{j} = \cos \theta \vec{e}_\phi, &
            \devp{\vec{e}_\phi}{\theta} &= 0, \\
            \devp{\vec{e}_\phi}{\phi} &= -\cos \phi \vec{i} - \sin \phi \vec{j} = - \cos \theta \vec{e}_\theta - \sin \theta \vec{e}_r.
        \end{align*}
    \end{solution}
  
    \question Seja o campo vetorial
    \[
    \vec{V} = \frac{y z}{r(x^2 + y^2)} \vec{i} - \frac{x z}{r (x^2 + y^2)} \vec{j},
    \]
    onde $r = \sqrt{x^2 + y^2 + z^2}$.
    \begin{parts}
        \part Usando coordenadas cartesianas, mostre que
        \[
        \nabla \times \vec{V} = \frac{\vec{r}}{r^3}.
        \]
        \begin{solution}
            Utilizando \eqref{eq:rot} temos
            \begin{align*}
                \nabla \times \vec{V} &= \begin{vmatrix}
                    \vec{i} & \vec{j} & \vec{k} \\
                    \devp{}{x} & \devp{}{y} & \devp{}{z} \\
                    \frac{yz}{r\left(x^2 + y^2\right)} & -\frac{xz}{r\left(x^2 + y^2\right)} & 0
                \end{vmatrix} \\
                &= \devp{\frac{yz}{r\left(x^2 + y^2\right)}}{z} \vec{j} + \devp{\frac{xz}{r\left(x^2 + y^2\right)}}{x} \vec{k} - \devp{\frac{yz}{r\left(x^2 + y^2\right)}}{y} \vec{k} - \devp{\frac{xz}{r\left(x^2 + y^2\right)}}{z} \vec{i} \\
                %&= \left(\frac{yr\left(x^2 + y^2\right) - yz\frac{1}{2}\left(r^{-1/2}2z\left(x^2 + y^2\right)\right)}{r^2 \left(x^2 + y^2\right)^2}\right)
                &= \left(\frac{-xz^2}{r^3 \left(x^2 + y^2\right)} + \frac{x}{r\left(x^2 + y^2\right)}\right) \vec{i} + \left(\frac{-yz^2}{r^3 \left(x^2 + y^2\right)} + \frac{y}{r \left(x^2 + y^2\right)}\right) \vec{j} \\ &\quad + \left(xz \left(\frac{-2x}{r\left(x^2 + y^2\right)^3} - \frac{x}{r^3 \left(x^2 + y^2\right)}\right) + \frac{z}{r}\frac{1}{x^2 + y^2}\right. \\ &\quad \left.- \left(yz \left(\frac{-2y}{r\left(x^2 + y^2\right)^3} - \frac{y}{r^3 \left(x^2 + y^2\right)}\right) + \frac{z}{r} \frac{1}{x^2 + y^2}\right)\right) \vec{k} \\
                &= \frac{1}{r^3} \left( \left(\frac{-xz^2 + x \left(x^2 + y^2 + z^2\right)}{x^2 + y^2}\right) \vec{i} + \left(\frac{-yz^2 + y\left( x^2 + y^2 + z^2 \right)}{x^2 + y^2}\right) \vec{j} \right. \\ & \quad \left. + \left(\frac{-2x^2 z \left(x^2 + y^2 + z^2\right) - x \left( x^2 + y^2 \right)}{\left(x^2 + y^2\right)^3} - \left( \frac{2 y^2 z \left( x^2 + y^2 + z^2 \right) - y \left( x^2 + y^2 \right)^2}{\left(x^2 + y^2\right)^3} \right)\right) \vec{k}\right) \\
                &= \frac{1}{r^3} \left( x \vec{i} + y \vec{j} + z \vec{k} \right) \\
                &= \frac{\vec{r}}{r^3}.
            \end{align*}
        \end{solution}
  
        \part Mostre que em termos de coordenadas esf\'{e}ricas
        \[
        \vec{V} = - \frac{\cot \theta}{r} \vec{e}_\phi.
        \]
        \begin{solution}
            Utilizando \eqref{eq:coor_esf}, escrevemos $r = x \vec{i} + y \vec{j} + z \vec{k}$ em coordenadas esf\'{e}ricas:
            \[
            \vec{r} = r \sin \theta \cos \phi \vec{i} + r \sin \theta \sin \phi \vec{j} + r \cos \theta \vec{k}.
            \]
  
            Calculando o fator de escala utilizando \eqref{eq:fator_escala} temos
            \begin{align*}
                h_r &= \sqrt{\sin^2 \theta \cos^2 \phi + \sin^2 \theta \sin^2 \phi + \cos^2 \theta} = 1, \\
                h_\theta &= \sqrt{r^2 \cos^2 \theta \cos^2 \phi + r^2 \cos^2 \theta \sin^2 \phi + r^2 \sin^2 \theta} = r, \\
                h_\phi &= \sqrt{r^2 \sin^2 \theta \sin^2 \phi + r^2 \sin^2 \theta \cos^2 \phi} = r \sin \theta.
            \end{align*}
            E calculando os vetores tangentes unit\'{a}rios utilizando \eqref{eq:vetor_tang_unit} temos
            \begin{align*}
                \vec{e}_r &= \sin \theta \cos \phi \vec{i} + \sin \theta \sin \phi \vec{j} + \cos \vec{k}, \\
                \vec{e}_\theta &= \cos \theta \cos \phi \vec{i} + \cos \theta \sin \phi \vec{j} - \sin \theta \vec{k}, \\
                \vec{e}_\phi &= -\sin \phi \vec{i} + \cos \phi \vec{j}.
            \end{align*}
  
            Agora, por meio de \eqref{eq:vetor_mud_coor}, escrevemos $\vec{i}$, $\vec{j}$ e $\vec{k}$ em fun\c{c}\~{a}o de $\vec{e}_r$, $\vec{e}_\theta$ e $\vec{e}_\phi$:
            \begin{align*}
                \vec{i} &= \sin \theta \cos \phi \vec{e}_r + \cos \theta \cos \phi \vec{e}_\theta - \sin \vec{e}_\phi, \\
                \vec{j} &= \sin \theta \sin \phi \vec{e}_r + \cos \theta \sin \phi \vec{e}_\theta + \cos \phi \vec{e}_\phi, \\
                \vec{k} &= \cos \theta \vec{e}_r - \sin \theta \vec{e}_\theta.
            \end{align*}
  
            Por fim, substituimos $x$, $y$, $z$, $\vec{i}$, $\vec{j}$ e $\vec{k}$ em $\vec{V}$:
            \begin{align*}
                \vec{V} &= \frac{\left( r \sin \theta \sin \phi \right)r \cos \theta}{r^3 \sin \theta} \left( \sin \theta \cos \phi \vec{e}_r + \cos \theta \cos \phi \vec{e}_\theta - \sin \phi \vec{e}_\phi \right) \\ & \quad - \frac{\left( r \sin \theta \cos \phi \right) r \cos \theta \left( \sin \theta \sin \phi \vec{e}_r + \cos \theta \sin \phi \vec{e}_\theta + \cos \phi \vec{e}_\phi \right)}{r^3 \sin \theta} \\
                &= \frac{1}{r \sin \theta} \left( \vec{e}_r \left( \sin \phi \cos \theta \sin \theta \cos \phi - \sin \theta \cos \phi \sin \phi \cos \theta \right) \right. \\ & \quad \left. + \vec{e}_\theta \left( \sin \phi \cos^2 \theta \cos \phi - \cos^2 \theta \sin \phi \cos \phi \right) + \vec{e}_\phi \left( - \sin^2 \phi \cos \theta - \cos^2 \phi \cos \theta \right) \right) \\
                &= - \frac{\cot \theta}{r} \vec{e}_\phi.
            \end{align*}
        \end{solution}
  
        \part Obtenha o resultado acima para $\nabla \times \vec{V}$ usando coordenadas esf\'{e}ricas.
        \begin{solution}
            Utilizando \eqref{eq:rot} temos que
            \begin{align*}
                \rot \vec{V} &= \frac{1}{r^2 \sin \theta} \begin{bmatrix}
                    \vec{e}_r & r \vec{e}_\theta & r \sin \theta \vec{e}_\phi \\
                    \devp{}{r} & \devp{}{\theta} & \devp{}{\phi} \\
                    0 & 0 & r \sin \theta \frac{-\cot \theta}{r}
                \end{bmatrix} \\
                &= \frac{1}{r^2 \sin \theta} \left( \devp{-\cos \theta}{\theta} \vec{e}_r + \devp{\cos \theta}{r} r \vec{e}_\theta \right) \\
                &= \frac{\vec{e}_r}{r^2} \\
                &= \frac{\vec{v}}{r^3}.
            \end{align*}
        \end{solution}
    \end{parts}
  
    \question Seja $f(r)$ uma fun\c{c}\~{a}o de $r = \sqrt{x^2 + y^2 + z^2}$. Mostre que
    \begin{align*}
        \nabla f(r) &= f'(r) \vec{e}_r, \\
        \nabla \cdot (\vec{e}_r f(r)) &= \frac{2}{r} f(r) + f'(r), \\
        \nabla \times (\vec{e}_r f(r)) &= 0
    \end{align*}
    usando:
    \begin{parts}
        \part coordenadas cartesianas,
        \begin{solution}
            % TODO Escrever a solu\c{c}\~{a}o.
        \end{solution}
  
        \part coordenads esf\'{e}ricas.
        \begin{solution}
            % TODO Escrever a solu\c{c}\~{a}o.
        \end{solution}
    \end{parts}
  
    \question Sejam
    \begin{align*}
        \nabla &= \vec{e}_r \devp{}{r} + \vec{e}_\theta \frac{1}{r} \devp{}{\theta} + \vec{e}_\phi \frac{1}{r \sin \theta} \devp{}{\phi}, \\
        \vec{V} &= V_r \vec{e}_r + V_\theta \vec{e}_\theta + V_\phi \vec{e}_\phi,
    \end{align*}
    e $\vec{V} \cdot \nabla$ o operador dado por
    \[
    \vec{V} \cdot \nabla = V_r \devp{}{r} + \frac{V_\theta}{r} \devp{}{\theta} + \frac{V_\phi}{r \sin \theta} \devp{}{\phi}.
    \]
    Mostre que $(\vec{V} \cdot \nabla) r = \vec{V}$.
    \begin{solution}
        Utilizando \eqref{eq:coor_esf}, escrevemos $r = x \vec{i} + y \vec{j} + z \vec{k}$ em coordenadas esf\'{e}ricas: 
        \[
        \vec{r} = r \sin \theta \cos \phi \vec{i} + r \sin \theta \sin \phi \vec{j} + r \cos \theta \vec{k}.
        \]
        
        Calculando o fator de escala utilizando \eqref{eq:fator_escala} temos
        \begin{align*}
            h_r &= \sqrt{\sin^2 \theta \cos^2 \phi + \sin^2 \theta \sin^2 \phi + \cos^2 \theta} = 1, \\
            h_\theta &= \sqrt{r^2 \cos^2 \theta \cos^2 \phi + r^2 \cos^2 \theta \sin^2 \phi + r^2 \sin^2 \theta} = r, \\
            h_\phi &= \sqrt{r^2 \sin^2 \theta \sin^2 \phi + r^2 \sin^2 \theta \cos^2 \phi} = r \sin \theta.
        \end{align*}
        E calculando os vetores tangentes unit\'{a}rios utilizando \eqref{eq:vetor_tang_unit} temos
        \begin{align*}
            \vec{e}_r &= \sin \theta \cos \phi \vec{i} + \sin \theta \sin \phi \vec{j} + \cos \vec{k}, \\
            \vec{e}_\theta &= \cos \theta \cos \phi \vec{i} + \cos \theta \sin \phi \vec{j} - \sin \theta \vec{k}, \\
            \vec{e}_\phi &= -\sin \phi \vec{i} + \cos \phi \vec{j}.
        \end{align*}
        
        Agora, por meio de \eqref{eq:vetor_mud_coor}, escrevemos $\vec{i}$, $\vec{j}$ e $\vec{k}$ em fun\c{c}\~{a}o de $\vec{e}_r$, $\vec{e}_\theta$ e $\vec{e}_\phi$:
        \begin{align*}
            \vec{i} &= \sin \theta \cos \phi \vec{e}_r + \cos \theta \cos \phi \vec{e}_\theta - \sin \vec{e}_\phi, \\
            \vec{j} &= \sin \theta \sin \phi \vec{e}_r + \cos \theta \sin \phi \vec{e}_\theta + \cos \phi \vec{e}_\phi, \\
            \vec{k} &= \cos \theta \vec{e}_r - \sin \theta \vec{e}_\theta.
        \end{align*}

        Substituindo $x$, $y$, $z$, $\vec{i}$, $\vec{j}$ e $\vec{k}$ em $\vec{r}$ temos
        \begin{align*}
            \vec{r} &= r \sin^2 \theta \cos^2 \phi \vec{e}_r + r \sin \theta \cos^2 \phi \cos \theta \vec{e}_\theta - r \sin^2 \theta \cos \phi \vec{e}_\phi \\
            & \quad + r \sin^2 \theta \sin^2 \phi \vec{e}_r + r \sin \theta \sin^2 \phi \cos \theta \vec{e}_\theta \\
            & \quad + r \sin \theta \cos \phi \sin \phi \vec{e}_\phi + r \cos^2 \theta \vec{e}_r - r \sin \theta \cos \theta \vec{e}_\theta \\
            &= \left( r \sin^2 \theta + r \cos^2 \theta \right) \vec{e}_r + \left( r \sin \theta \cos \theta - r \sin \theta \cos \theta \right) \vec{e}_\theta \\
            &= r \vec{e}_r.
        \end{align*}

        Por fim,
        \begin{align*}
            \left( \vec{V} \cdot \nabla \right) \vec{r} &= \left( V_r \devp{}{r} + \frac{V_\theta}{r} \devp{}{\theta} + \frac{V_\phi}{r \sin \theta} \devp{}{\phi} \right) r \vec{e}_r \\
            &= V_r \vec{e}_r + \frac{V_\theta}{r} \devp{r \vec{e}_r}{\theta} + \frac{V_\phi}{r \sin \theta} \devp{r \vec{e}_r}{\phi} \\
            &= V_r \vec{e}_r + V_\theta \vec{e}_\theta + \frac{V_\phi \sin \theta \vec{e}_\phi}{\sin \theta} \\
            &= V_r \vec{e}_r + V_\theta \vec{e}_\theta + V_\phi \vec{e}_\phi.
        \end{align*}
    \end{solution}

    % TODO Incluir T1 de 2008.

    % TODO Incluir T2 de 2008.

    \question[P1 de 2008] Considere o campo vetorial
    \begin{align*}
        \vec{F} &= (x^2 + y^2 + z^2)^n (x \vec{x} + y \vec{y} + z \vec{z}).
    \end{align*}
    \begin{parts}
        \part Escreva $\vec{F}$ em coordenadas cil\'{i}ndricas.
        \begin{solution}
            Em coordenadas cil\'{i}ndricas temos
            \begin{align*}
                \vec{F} &= (r^2 + z^2)^n (r \vec{r} + z \vec{z}).
            \end{align*}
        \end{solution}

        \part Calcule $\nabla \cdot \vec{F}$.
        \begin{solution}
            Temos que
            \begin{align*}
                \nabla \cdot \vec{F} &= \devp{x (x^2 + y^2 + z^2)^n}{x} + \devp{y (x^2 + y^2 + z^2)^n}{y} + \devp{z (x^2 + y^2 + z^2)^n}{z} \\  \\
                &= 3 (x^2 + y^2 + z^2) + \left( x \devp{}{x} + y \devp{}{x} + z \devp{}{z} \right) \left( x^2 + y^2 + z^2 \right)^n \\
                &= (2n + 3) \left( x^2 + y^2 + z^2 \right)^n.
            \end{align*}
        \end{solution}

        \part Calcule $\nabla \times \vec{F}$.
        \begin{solution}
            Temos que, para a componente $\vec{i}$,
            \begin{align*}
                \left( \devp{F_z}{y} - \devp{F_y}{z} \right) \vec{i} &= z \devp{(x^2 + y^2 + z^2)^n}{y} - y \devp{(x^2 + y^2 + z^2)^n}{z} \\
                &= n (x^2 + y^2 + z^2)^{n - 1} \left( z \devp{y^2}{y} - y \devp{z^2}{z} \right) \\
                &= 0.
            \end{align*}
            Portanto, $\nabla \times \vec{F} = 0$.
        \end{solution}

        \part Determine uma fun\c{c}\~{a}o escalar $\phi$ tal que $\vec{F} = - \nabla \phi$.
        \begin{solution}
            Em coordenadas esf\'{e}ricas temos que $\vec{F} = r^{2n + 1} \vec{r}$ e portanto $\phi = - r^{2n + 2} / (2n + 2)$, $n \neq - 1$, e $\phi = - \ln r$, $n = -1$.
        \end{solution}

        \part Para que valores de $n$ a fun\c{c}\~{a}o escalar $\phi$ diverge na origem?
        \begin{solution}
            A fun\c{c}\~{a}o $\phi$ diverge na origem para $n < 0$.
        \end{solution}
    \end{parts}

    \question[E de 2008] Considere o seguinte campo vetorial em $\mathbb{R}^3$
    \begin{align*}
        \vec{F} &= \ln\left( x^2 + y^2 + z^2 \right) \left( x \vec{x} + y \vec{y} + z \vec{z} \right).
    \end{align*}
    \begin{parts}
        \part Escreva $\vec{F}$ em coordenadas cil\'{i}ndricas.
        \begin{solution}
            % TODO Detalhar solu\c{c}\~{a}o.
            Temos que
            \begin{align*}
                \vec{F} &= \ln(r^2 + z^2) \left( r \vec{r} + z \vec{z} \right).
            \end{align*}
        \end{solution}

        \part Calcule $\oint_\Gamma \vec{F} \cdot \mathrm{d}\vec{s}$, sendo $\Gamma$ o c\'{i}rculo unit\'{a}rio com centro na origem contido no plano $z = 0$.
        \begin{solution}
            % TODO Detalhar solu\c{c}\~{a}o.
            A integral \'{e} zero pois o campo \'{e} radial.
        \end{solution}

        \part Calcule $\nabla \cdot \vec{F}$.
        \begin{solution}
            % TODO Detalhar solu\c{c}\~{a}o.
            Temos que
            \begin{align*}
                \sum_i \left[ \partial_i \left( i \ln(x^2 + y^2 + z^2) \right) \right] &= 3 \ln(x^2 + y^2 + z^2) + 2
            \end{align*}
            e portanto $\nabla \cdot \vec{F} = 0$.
        \end{solution}

        \part Determine a fun\c{c}\~{a}o escalar $\phi$ tal que $\vec{F} = - \nabla \phi$.
        \begin{solution}
            Em coordenadas esf\'{e}ricas temos que $\vec{F} = 2 r \vec{r} \ln r$, de onde tem-se que
            \begin{align*}
                \phi = -2 \int r \ln r \,\mathrm{d}r \\
                &= r~2 \left( 1/2 - \ln r \right).
            \end{align*}
        \end{solution}
    \end{parts}
    
    \question[E de 2010] Considere o campo vetorial dado, em coordenadas esf\'{e}ricas, por
    \begin{align*}
        \vec{F} &= \frac{2 P \cos \theta}{r^3} \vec{r} + \frac{P \sin \theta}{r^3} \vec{\theta}, r \geq P/2,
    \end{align*}
    onde $P > 0$. Calcule $\int_C \vec{F} \cdot \mathrm{d}\vec{r}$ onde $C$ \'{e} o segmento de reta orientado que liga $A$ a $B$, onde $A$ e $B$ s\~{a}o dados por $A = P(0,0,-1)$ e $B = P(1,1,1)$ em coordenadas cartesianas.
    \begin{solution}
        % TODO Escrever solu\c{c}\~{a}o.
    \end{solution}

    \question[T1 de 2011] Sejam $(u, v, z)$ as coordenadas cil\'{i}ndricas parab\'{o}licas, definidas como
    \begin{align*}
        x &= (1/2)(u^2 - v^2), & y &= uv, & z &= z,
    \end{align*}
    com $-\infty < u < +\infty$, $v \geq 0$, $-\infty < z < +\infty$.
    \begin{parts}
        \part Mostre que esse sistema de coordenadas \'{e} ortogonal.
        \begin{solution}
           Temos que
           \begin{align*}
               r &= (1/2)(u^2 - v^2) \vec{i} + u v \vec{j} + z \vec{k}
           \end{align*}
           e portanto
           \begin{align*}
               \devp{\vec{r}}{u} &= u\vec{i} + v \vec{j}, \\
               \devp{\vec{r}}{v} &= -v \vec{i} + v \vec{j}, \\
               \devp{\vec{r}}{v} & = k.
           \end{align*}
           Logo,
           \begin{align*}
               h_u &= \sqrt{u^2 + v^2}, \\
               h_v &= \sqrt{u^2 + v^2}, \\
               h_z &= 1
           \end{align*}
           e
           \begin{align*}
               \vec{e}_u &= (u\vec{i} + v\vec{j})/ \sqrt{u^2 + v^2}, \\
               \vec{e}_v &= (-v\vec{i} + u\vec{j})/ \sqrt{u^2 + v^2}, \\
               \vec{e}_z &= \vec{k}.
           \end{align*}
           Ent\~{a}o
           \begin{align*}
               \vec{e}_u \vec{e}_v &= -uv + vu = 0, \\
               \vec{e}_u \vec{e}_z &= 0, \\
               \vec{e}_v \vec{e}_z &= 0
           \end{align*}
           e assim concluimos que o sistema de coordenadas \'{e} ortogonal.
        \end{solution}

        \part Seja $r$ o vetor posi\c{c}\~{a}o, $r = x\vec{i} + y\vec{j} + z\vec{k}$. Mostre que em coordenadas cil\'{i}ndricas parab\'{o}licas temos
        \begin{align*}
            r = (1/2) \sqrt{u^2 + v^2} (u\vec{e}_u + v\vec{e}_v) + z\vec{e}_z.
        \end{align*}
        \begin{solution}
            Utilizando i\eqref{eq:vetor_mud_coor} temos que
            \begin{align*}
                \vec{i} &= \left( \vec{i} \cdot \vec{e}_u \right) \vec{e}_u + \left( \vec{i} \cdot \vec{e}_v \right) \vec{e}v + \left( \vec{i} \cdot \vec{e}_z \right) \vec{e}_z \\
                &= \frac{u}{\sqrt{u^2 + v^2}} \vec{e}_u - \frac{u}{\sqrt{u^2 + v^2}} \vec{e}_v, \\
                \vec{j} &= \left( \vec{j} \cdot \vec{e}_u \right) \vec{e}_u + \left( \vec{j} \cdot \vec{e}_v \right) \vec{e}_v + \left( \vec{j} + \vec{e}_z \right) \vec{e}_z \\
                &= \frac{v}{\sqrt{u^2 + v^2}} \vec{e}_u + \frac{u}{\sqrt{u^2 + v^2}} \vec{e}_v, \\
                \vec{k} &= \vec{e}_z.
            \end{align*}
            Portanto,
            \begin{align*}
                \begin{split}
                    \vec{r} &= \left( (u^2 - v^2) / 2 \right) \left( \frac{u}{\sqrt{u^2 + v^2}} \vec{e}_u - \frac{v}{\sqrt{u^2 + v^2}} \vec{e}_v \right) \\
                    &\quad {}+ uv \left( \frac{v}{\sqrt{u^2 + v^2}} \vec{e}_u + \frac{u}{\sqrt{u^2 + v^2}} \vec{e}_v \right) + z \vec{e}_z
                \end{split} \\
                &= \frac{u (u^2 - v^2}{2 \sqrt{u^2 + v^2}} \vec{e}_u - \frac{v (u^2 - v^2}{2 \sqrt{u^2 + v^2}} \vec{e}_v + \frac{u v^2}{\sqrt{u^2 + v^2}} \vec{e}_u + \frac{u^2 v}{\sqrt{u^2 + v^2}} \vec{e}_v + z \vec{e}_z \\
                &= \frac{u (u^2 + v^2)}{2 \sqrt{u^2 + v^2}} \vec{e}_u + \frac{u (u^2 + v^2)}{2 \sqrt{u^2 + v^2}} \vec{e}_v + z \vec{e}_z \\
                &= \frac{\sqrt{u^2 + v^2}}{2} (u \vec{e}_u + v \vec{e}_v) + z \vec{e}_z.
            \end{align*}
        \end{solution}

        \part Usando as coordenas, e sabendo que em um sistema de coordenadas curvil\'{i}neas ortogonal temos
        \begin{align*}
            \nabla \cdot \vec{V} &= \frac{1}{h_1 h_2 h_ 3} \left[ \devp{}{q_1}(h_2 h_3 V_1) + \devp{}{q_2}(h_3 h_1 V_2) + \devp{}{q_3}(h_1 h_2 V_3) \right],
        \end{align*}
        mostre que $\nabla \cdot \vec{r} = 3$.
        \begin{solution}
            Pelo item anterior temos que $r_u = (1/2) u \sqrt{u^2 + v^2}$, $r_v = (1/2) v \sqrt{u^2 + v^2}$ e $r_z = z$. Logo,
            \begin{align*}
                \nabla \cdot \vec{r} &= \frac{1}{u^2 + v^2} \left[ \devp{}{u}\left( \frac{(u^2 + v^2)u}{2} \right) + \devp{}{v}\left( \frac{(u^2 + v^2)v}{2} + \devp{}{z}\left( (u^2 + v^2) z \right) \right) \right] \\
                &= \frac{1}{u^2 + v^2} \left[ (3/2)u^2 + (1/2)v^2 + (1/2)u^2 + (3/2)v^2 + u^2 + v^2 \right] \\
                &= \frac{1}{u^2 + v^2} \frac{6}{2} (v^2 + u^2) = 3.
            \end{align*}
        \end{solution}
    \end{parts}

    \question[P1 de 2011]
    \begin{parts}
        \part Sejam as coordenadas esf\'{e}ricas $(r, \theta, \phi)$ dadas por
        \begin{align*}
            x &= r \sin \theta \cos \phi, & y &= r \sin \theta \sin \phi, & z &= r \cos \theta,
        \end{align*}
        onde $0 \leq r \leq \infty$, $0 \leq \theta \leq \pi$ e $0 \leq \phi \leq 2 \pi$. Mostre que os vetores tangentes unit\'{a}rios $\left\{ \vec{e}_r, \vec{e}_\theta, \vec{e}_\phi \right\}$ s\~{a}o dados por
        \begin{align*}
            \vec{e}_r &= \sin \theta \cos \phi \vec{i} + \sin \theta \sin \phi \vec{j} + \cos \theta \vec{k}, \\
            \vec{e}_\theta &= \cos \theta \cos \phi \vec{i} + \cos \theta \sin \phi \vec{j} - \sin \theta \vec{k}, \\
            \vec{e}_\phi &= - \sin \phi \vec{i} + \cos \phi \vec{j}.
        \end{align*}
        onde $\left\{ \vec{i}, \vec{j}, \vec{k} \right\}$ s\~{a}o os vetores unit\'{a}rios cartesianos tais que $\vec{r} = x \vec{i} + y \vec{j} + z \vec{k}$.
        \begin{solution}
            Temos que $\vec{r} = r \sin \theta \cos \phi \vec{i} + r \sin \theta \cos \phi \vec{j} + r \cos \theta \vec{k}$ ent\~{a}o utilizando \eqref{eq:vetor_tangente} obtemos
            \begin{align*}
                \devp{\vec{r}}{r} &= \sin \theta \cos \phi \vec{i} + \sin \theta \sin \phi \vec{j} + \cos \theta \vec{k}, \\
                \devp{\vec{r}}{\theta} &= r \cos \theta \cos \phi \vec{i} + r \cos \theta \sin \phi \vec{j} - r \sin \theta \vec{k}, \\
                \devp{\vec{r}}{\phi} &= -r \sin \theta \sin \phi \vec{i} + r \sin \theta \cos \phi \vec{j}.
            \end{align*}
            Utilizando \eqref{eq:fator_escala} obtemos
            \begin{align*}
                h_r^2 &= \sin^2 \theta \cos^2 \phi + \sin^2 \theta \sin^2 \phi + \cos^2 \theta = 1, \\
                h_\theta^2 &= r^2 \cos^2 \theta \cos^2 \phi + r^2 \cos^2 \theta \sin^2 \phi + r^2 \sin^2 \theta = r^2, \\
                h_\phi^2 &= r^2 \sin^2 \theta \sin^2 \phi + r^2 \sin^2 \theta \cos^2 \phi = r^2 \sin^2 \theta.
            \end{align*}
            Utilizando \eqref{eq:vetor_tang_unit} obtemos
            \begin{align*}
                \vec{e}_r &= \sin \theta \cos \phi \vec{i} + \sin \theta \cos \phi \vec{j} + \cos \theta \vec{k}, \\
                \vec{e}_\theta &= \cos \theta \cos \phi \vec{i} + \cos \theta \sin \phi \vec{j} - \sin \theta \vec{k}, \\
                \vec{e}_\phi &= - \sin \phi \vec{i} + \cos \phi \vec{j}.
            \end{align*}
        \end{solution}

        \part Sejam $\nabla$ e $\vec{V}$ dados por
        \begin{align*}
            \nabla &= \vec{e}_r \devp{}{r} + \vec{e}_\theta \frac{1}{r} \devp{}{\theta} + \vec{e}_\phi \frac{1}{r \sin \theta} \devp{}{\phi}, & \vec{V} &= V_r \vec{e}_r + V_\theta \vec{e}_\theta + V_\phi \vec{e}_\phi,
        \end{align*}
        e $\vec{V} \cdot \nabla$ o operador dado por
        \begin{align*}
            \vec{V} \cdot \nabla &= V_r \devp{}{r} + \frac{V_\theta}{r} \devp{}{\theta} + \frac{V_\phi}{r \sin \theta} \devp{}{\phi}.
        \end{align*}
        Mostre que $(\vec{V} \cdot \nabla) \vec{r} = \vec{V}$.
        \begin{solution}
            Tomemos que $\vec{r} = r \vec{e}_r$ e 
            \begin{align*}
                (\vec{V} \cdot \nabla) \vec{r} &= V_r \devp{r \vec{e}_r}{r} + \frac{V_\theta}{r} \devp{r \vec{e}_r}{\theta} + \frac{V\phi}{r \sin \theta} \devp{r \vec{e}_r}{\phi}.
            \end{align*}
            Mas
            \begin{align*}
                \devp{r \vec{e}_r}{r} &= \left( \devp{r}{r} \right) \vec{e}_r + r \left( \devp{\vec{e}_r}{r} \right) \\
                &= \vec{e}_r, \\
                \devp{r \vec{e}_r}{\theta} &= \left( \devp{r}{\theta} \right) \vec{e}_r + r \left( \devp{\vec{e}_r}{\theta} \right) \\
                &= r \left( \cos \theta \cos \phi \vec{i} + \cos \theta \sin \phi \vec{j} - \sin \theta \vec{k} \right) = r \vec{e}_\theta, \\
                \devp{r \vec{e}_r}{\phi} &= \left( \devp{r}{\phi} \right) \vec{e}_r + r \left( \devp{\vec{e}_r}{\phi} \right) \\
                &= r \left( - \sin \theta \sin \phi \vec{i} + \sin \theta \cos \phi \vec{j} \right) = r \sin \theta \vec{e}_\phi.
            \end{align*}
            Logo,
            \begin{align*}
                \left( \vec{V} \cdot \nabla \right) \vec{r} &= V_r \vec{e}_r + \frac{V_\theta}{r} r \vec{e}_\theta + \frac{V_\phi}{r \sin \theta} r \sin \theta \vec{e}_\phi \\
                &= V_r \vec{e}_r + V_\theta \vec{e}_\theta + V_\phi \vec{e}_\phi = \vec{V}.
            \end{align*}
        \end{solution}
    \end{parts}

    \question[T1 de 2012] Sejam as coordenads esf\'{e}ricas $(r, \theta, \phi)$ dadas por
    \begin{align*}
        x &= r \sin \theta \cos \phi, & y &= r \sin \theta \cos \phi, & z &= r \cos \theta,
    \end{align*}
    onde $0 \leq r \leq \infty$, $0 \leq \theta \leq \pi$ e $0 \leq \phi \leq 2 \pi$.
    \begin{parts}
        \part Mostre que os fatores de escala s\~{a}o dados por
        \begin{align*}
            h_r &= 1, & h_\theta &= r, & h_\phi &= r \sin \theta.
        \end{align*}
        \begin{solution}
            Utilizando \eqref{eq:fator_escala} temos que
            \begin{align*}
                h_r^2 &= (\sin \theta \cos \phi)^2 + (\sin \theta \cos \phi)^2 + \cos^2 \theta \\
                &= \sin^2 \theta + \cos^2 \theta = 1, \\
                h_\theta^2 &= (r \cos \theta \cos \phi)^2 + (r \cos \theta \sin \phi)^2 + (-r \sin \theta)^2 \\
                &= r^2 \cos^2 \theta + r^2 \sin \theta = r^2, \\
                h_\phi^2 &= (-r \sin \theta \sin \phi)^2 + (r \sin \theta \cos \phi)^2 \\
                &= r^2 \sin^2 \theta.
            \end{align*}
        \end{solution}

        \part Seja $\vec{A}$ o campo vetorial dado por
        \begin{align*}
            \vec{A} &= (r/3) \sin \theta \vec{e}_\phi.
        \end{align*}
        Usando coordenadas esf\'{e}ricas, calcule $\vec{B} = \nabla \times \vec{A}$ e $\rho = \nabla \cdot \vec{B}$.
        \begin{solution}
            Se $\vec{A} = (r/3) \sin \theta \vec{e}_\phi$ ent\~{a}o $A_r = 0$, $A_\theta = 0$ e $A_\phi = (r/3) \sin \theta$. Portanto
            \begin{align*}
                \nabla \times \vec{A} &= \frac{1}{r^2 \sin \theta} \begin{vmatrix}
                    \vec{e}_r & r \vec{e}_\theta & r \sin \theta \vec{e}_\phi \\
                    \devp{}{r} & \devp{}{\theta} & \devp{}{\phi} \\
                    0 & 0 & (r^2 \sin^2 \theta) / 3
                \end{vmatrix} \\
                &= \frac{1}{r^2 \sin \theta} \left( \vec{e}_r (2r^2 / 3) \sin \theta \cos \theta - r \vec{e}_\theta (2r/3) \sin^2 \theta \right) \\
                &= (2/3) \cos \theta \vec{e}_r - (2/3) \sin \theta \vec{e}_\theta.
            \end{align*}
            Ent\~{a}o $\vec{B} = (2/3) \cos \theta \vec{e}_r - (2/3) \sin \theta \vec{e}_\theta$ e portanto $B_r = (2/3) \cos \theta$, $B_\theta = (-2/3) \sin \theta$ e $B_\phi = 0$. Logo,
            \begin{align*}
                \nabla \cdot \vec{B} &= \frac{1}{r^2 \sin \theta} \left[ \devp{}{r}(r^2 \sin \theta (2/3) \cos \theta) + \devp{}{\theta} (r \sin \theta (-2/3) \sin \theta) \right] \\
                &= \frac{1}{r^2 \sin \theta} \left[ (4/3) r \sin \theta \cos \theta - (4/3) r \sin \theta \cos \theta \right] = 0.
            \end{align*}
        \end{solution}
    \end{parts}

    \question[P1 de 2012] Sejam as coordenadas esf\'{e}ricas $(r, \theta, \phi)$ dadas por
    \begin{align*}
        x &= r \sin \theta \cos \phi, & y &= r \sin \theta \cos \phi, & z &= r \cos \theta,
    \end{align*}
    onde $0 \leq r \leq \infty$, $0 \leq \theta \leq \pi$ e $0 \leq \phi \leq 2 \pi$. Mostre que
    \begin{align*}
        \nabla \times (\cos \theta \nabla \phi) = \nabla (1/r).
    \end{align*}
    \begin{solution}
        Temos que
        \begin{align*}
            \nabla \times (\cos \theta \nabla \phi) &= \nabla(\cos \theta) \times \nabla \phi + \cos \theta \underbrace{\nabla \times (\nabla \phi)}_{= 0}.
        \end{align*}
        Como
        \begin{align*}
            \nabla(\cos \theta) &= \frac{1}{h_r} \underbrace{\devp{\cos \theta}{r}}_{= 0} \vec{e}_r + \frac{1}{h_\theta} \devp{\cos \theta}{\theta} \vec{e}_\theta + \frac{1}{h_\phi} \underbrace{\devp{\cos \theta}{\phi}}_{= 0} \vec{e}_\phi \\
            &= \frac{1}{r} (- \sin \theta) \vec{e}_\theta, \\
            \nabla \phi &= \frac{1}{h_r} \underbrace{\devp{\phi}{r}}_{= 0} \vec{e}_r + \frac{1}{h_\theta} \underbrace{\devp{\phi}{\theta}}_{= 0} \vec{e}_\theta + \frac{1}{h_\phi} \devp{\phi}{\phi} \vec{e}_\phi \\
            &= \frac{1}{r \sin \theta} \vec{e}_\phi
        \end{align*}
        ent\~{a}o
        \begin{align*}
            \nabla \times (\cos \theta \nabla \phi) &= - \frac{\sin \theta}{r} \vec{e}_\theta \times \frac{1}{r \sin \theta} \vec{e}_\phi = - \frac{1}{r^2} \vec{e}_r.
        \end{align*}
        Por outro lado
        \begin{align*}
            \nabla\left( \frac{1}{r} \right) &= \frac{1}{h_r} \devp{}{r} \left( \frac{1}{r} \right) \vec{e}_r + \frac{1}{h_\theta} \underbrace{\devp{}{\theta} \left( \frac{1}{r} \right)}_{= 0} \vec{e}_\theta + \frac{1}{h_\phi} \underbrace{\devp{}{\phi} \left( \frac{1}{r} \right)}_{= 0} \vec{e}_\phi = - \frac{1}{r^2} \vec{e}_r.
        \end{align*}
        Portanto, concluimos que $\nabla \times (\cos \theta \nabla \phi) = \nabla(1/r)$.
    \end{solution}
\end{questions}
\end{document}


